\chapter{Analytical mechanics}
\label{ch:analytical}

\buildson{Newtonian mechanics (Chapter~\ref{ch:newton}), the harmonic oscillator (Chapter~\ref{ch:sho})}

\chapterdek{The same physics as Newton, reborn as a single principle --- nature is lazy --- and the doorway through which every modern theory enters.}

\begin{prototype}
A ball thrown from point $A$ at time $t_1$ to point $B$ at time $t_2$. Of all the paths it could conceivably take between those two events, it takes the one that makes a single number, the \emph{action}, stationary.
\end{prototype}

Newton's mechanics is about forces and vectors. It works, but it gets awkward fast: a bead on a curved wire, a double pendulum, or a spinning top all require you to find constraint forces you do not actually care about. In the eighteenth and nineteenth centuries Lagrange and Hamilton reformulated mechanics in terms of energies and a variational principle. The new formulation is easier to use, makes symmetries and conservation laws obvious, and --- this is the real reason it matters --- is the language in which quantum mechanics, field theory, and general relativity are written.

\section{The principle of stationary action}

Describe a system by any convenient set of \emph{generalized coordinates} $q_1,\dots,q_n$: angles, distances, whatever fits the problem. Define the \textbf{Lagrangian}
\[
  L(q,\dot q,t) = T - V,
\]
kinetic minus potential energy, and the \textbf{action} of a path $q(t)$ as
\[
  S[q] = \int_{t_1}^{t_2} L\big(q(t),\dot q(t),t\big)\,\dd t .
\]
The action is a \emph{functional}: feed it a whole path, and it returns a number.

\begin{law}[Hamilton's principle]
The path a system actually follows between fixed endpoints $q(t_1)$ and $q(t_2)$ is one for which the action is stationary: changing the path by a small amount $\delta q(t)$, vanishing at the endpoints, changes $S$ only at second order.
\end{law}

To turn this into an equation, vary the path, $q\to q + \delta q$, and expand:
\[
  \delta S = \int_{t_1}^{t_2}\left(\pdv{L}{q}\,\delta q + \pdv{L}{\dot q}\,\delta\dot q\right)\dd t
  = \int_{t_1}^{t_2}\left(\pdv{L}{q} - \dv{}{t}\pdv{L}{\dot q}\right)\delta q\,\dd t ,
\]
where we integrated the second term by parts and used $\delta q(t_1) = \delta q(t_2) = 0$. For this to vanish for \emph{every} $\delta q$, the bracket must vanish at every moment:

\begin{law}[The Euler--Lagrange equations]
\[
  \dv{}{t}\pdv{L}{\dot q_i} = \pdv{L}{q_i}, \qquad i = 1,\dots,n.
\]
\end{law}

\begin{example}[Recovering Newton]
For a particle in a potential, $L = \tfrac12 m\dot x^2 - V(x)$. Then $\pdv{L}{\dot x} = m\dot x$ and $\pdv{L}{x} = -V'(x)$, so the Euler--Lagrange equation is $m\ddot x = -V'(x)$. Newton's second law, recovered from a minimization principle.
\end{example}

\begin{example}[The pendulum, with no forces in sight]
With angle $\theta$ as the coordinate, $T=\tfrac12 m\ell^2\dot\theta^2$ and $V = -mg\ell\cos\theta$. The Euler--Lagrange equation gives $m\ell^2\ddot\theta = -mg\ell\sin\theta$ immediately. The tension in the string never appears, because the coordinate $\theta$ automatically respects the constraint. This is the practical superpower of Lagrangian mechanics: \emph{choose coordinates that respect the constraints, and the constraint forces disappear}.
\end{example}

\begin{example}[Polar coordinates, and a force that comes from nowhere]
A particle in a central potential $V(r)$, in the plane: $L = \tfrac12 m(\dot r^2 + r^2\dot\phi^2) - V(r)$. The $r$ equation is
\[
  m\ddot r = mr\dot\phi^2 - V'(r),
\]
and the term $mr\dot\phi^2$ is the centrifugal force, produced automatically by the kinetic energy in curvilinear coordinates. The $\phi$ equation says $\dv{}{t}(mr^2\dot\phi) = 0$: angular momentum is conserved, because $L$ does not depend on $\phi$.
\end{example}

\section{Momenta, symmetry, and Noether's theorem}

The quantity $p_i = \pdv{L}{\dot q_i}$ is the \emph{generalized momentum} conjugate to $q_i$. For an ordinary coordinate it is ordinary momentum; for an angle it is angular momentum. If $L$ happens not to depend on some $q_i$ (a \emph{cyclic} coordinate), the Euler--Lagrange equation says $\dot p_i = 0$: that momentum is conserved.

This is a special case of the most beautiful theorem in classical physics.

\begin{theorem}[Noether's theorem]
For every continuous symmetry of the action there is a conserved quantity.
\begin{itemize}[itemsep=2pt]
  \item Symmetry under translation in time $\Longrightarrow$ conservation of energy.
  \item Symmetry under translation in space $\Longrightarrow$ conservation of momentum.
  \item Symmetry under rotation $\Longrightarrow$ conservation of angular momentum.
\end{itemize}
\end{theorem}

Read that list again. Energy is conserved \emph{because} the laws of physics are the same today as they were yesterday. Momentum is conserved \emph{because} the laws are the same here as over there. These conservation laws are not separate facts about nature; they are shadows of its symmetries. When a symmetry is broken --- a ball rolling on a surface whose shape changes with time, say --- the corresponding quantity is no longer conserved.

The time-translation case deserves an explicit formula. If $L$ has no explicit time dependence, then
\[
  H = \sum_i p_i\dot q_i - L
\]
is conserved. For ordinary systems, $H = T+V$, the total energy.

\begin{keyidea}
Symmetries give conservation laws. In modern physics this runs in the other direction too: we \emph{postulate} symmetries and deduce the laws. Electric charge is conserved because of a symmetry of the quantum phase of the electron (Chapter~\ref{ch:qft}), and the forces of the Standard Model are dictated by symmetry (Chapter~\ref{ch:particle}).
\end{keyidea}

\section{Variational principles everywhere}

Hamilton's principle is one example of a pattern that runs through all of physics: \emph{the actual configuration extremizes something}.
\begin{itemize}[itemsep=3pt]
  \item \textbf{Fermat's principle} (Chapter~\ref{ch:optics}): light travels between two points along the path of stationary travel time. Snell's law of refraction falls out immediately.
  \item \textbf{The brachistochrone}: the curve of fastest descent between two points under gravity is not a straight line but a cycloid. Johann Bernoulli posed the problem in 1696; Newton reportedly solved it overnight.
  \item \textbf{Geodesics}: the shortest path between two points on a curved surface. In general relativity (Chapter~\ref{ch:gr}), freely falling objects follow geodesics of spacetime --- which is to say, they extremize their proper time.
  \item \textbf{Soap films} minimize their area, and a hanging chain minimizes its potential energy (the curve is a catenary, $y = a\cosh(x/a)$).
\end{itemize}

\section{Hamiltonian mechanics}

Lagrange's equations are $n$ second-order equations. Hamilton rewrote them as $2n$ first-order equations by treating the momenta $p_i$ as independent variables alongside the $q_i$. Trading $\dot q$ for $p$ by a \emph{Legendre transform} gives the \textbf{Hamiltonian} $H(q,p) = \sum p_i\dot q_i - L$, and the equations of motion become beautifully symmetric:

\begin{law}[Hamilton's equations]
\[
  \dot q_i = \pdv{H}{p_i}, \qquad \dot p_i = -\pdv{H}{q_i}.
\]
\end{law}

For the harmonic oscillator, $H = p^2/2m + \tfrac12 kx^2$, and Hamilton's equations give $\dot x = p/m$, $\dot p = -kx$: the familiar ellipse in phase space (Chapter~\ref{ch:sho}). In general, the state of a system is a single point in the $2n$-dimensional \emph{phase space}, and Hamilton's equations define a flow: every point moves along a definite trajectory, and trajectories never cross.

\begin{center}
\begin{minipage}{\linewidth}\centering
\begin{tikzpicture}
  \draw[-{Stealth[length=6pt]}] (-3.0,0) -- (3.15,0) node[right] {$\theta$};
  \draw[-{Stealth[length=6pt]}] (0,-2.55) -- (0,2.75) node[above] {$p$};
  \foreach \x/\l in {-2.5/{-\pi}, 2.5/{\pi}} \draw (\x,0.07) -- (\x,-0.07) node[below, font=\small] {$\l$};
  \draw[ascblue, thick] plot[smooth] coordinates {(-0.575,0.000) (-0.574,0.035) (-0.571,0.069) (-0.566,0.103) (-0.559,0.137) (-0.550,0.171) (-0.540,0.204) (-0.527,0.236) (-0.512,0.268) (-0.496,0.299) (-0.478,0.329) (-0.459,0.358) (-0.437,0.385) (-0.415,0.412) (-0.390,0.437) (-0.365,0.460) (-0.338,0.483) (-0.310,0.503) (-0.281,0.522) (-0.251,0.538) (-0.220,0.553) (-0.188,0.567) (-0.156,0.578) (-0.123,0.586) (-0.090,0.593) (-0.056,0.598) (-0.023,0.601) (0.011,0.601) (0.045,0.599) (0.079,0.595) (0.112,0.589) (0.145,0.581) (0.178,0.570) (0.210,0.558) (0.241,0.544) (0.271,0.527) (0.301,0.509) (0.329,0.489) (0.356,0.468) (0.382,0.445) (0.407,0.420) (0.430,0.394) (0.452,0.367) (0.472,0.339) (0.490,0.309) (0.507,0.278) (0.522,0.247) (0.536,0.215) (0.547,0.182) (0.556,0.149) (0.564,0.115) (0.570,0.081) (0.573,0.046) (0.575,0.000)} -- plot[smooth] coordinates {(0.575,-0.000) (0.574,-0.035) (0.571,-0.069) (0.566,-0.103) (0.559,-0.137) (0.550,-0.171) (0.540,-0.204) (0.527,-0.236) (0.512,-0.268) (0.496,-0.299) (0.478,-0.329) (0.459,-0.358) (0.437,-0.385) (0.415,-0.412) (0.390,-0.437) (0.365,-0.460) (0.338,-0.483) (0.310,-0.503) (0.281,-0.522) (0.251,-0.538) (0.220,-0.553) (0.188,-0.567) (0.156,-0.578) (0.123,-0.586) (0.090,-0.593) (0.056,-0.598) (0.023,-0.601) (-0.011,-0.601) (-0.045,-0.599) (-0.079,-0.595) (-0.112,-0.589) (-0.145,-0.581) (-0.178,-0.570) (-0.210,-0.558) (-0.241,-0.544) (-0.271,-0.527) (-0.301,-0.509) (-0.329,-0.489) (-0.356,-0.468) (-0.382,-0.445) (-0.407,-0.420) (-0.430,-0.394) (-0.452,-0.367) (-0.472,-0.339) (-0.490,-0.309) (-0.507,-0.278) (-0.522,-0.247) (-0.536,-0.215) (-0.547,-0.182) (-0.556,-0.149) (-0.564,-0.115) (-0.570,-0.081) (-0.573,-0.046) (-0.575,-0.000)} -- cycle;
  \draw[ascblue, thick] plot[smooth] coordinates {(-1.008,0.000) (-1.007,0.037) (-1.004,0.073) (-1.001,0.110) (-0.995,0.146) (-0.988,0.183) (-0.980,0.219) (-0.970,0.255) (-0.958,0.291) (-0.945,0.326) (-0.931,0.361) (-0.915,0.396) (-0.898,0.431) (-0.879,0.465) (-0.859,0.498) (-0.838,0.531) (-0.815,0.563) (-0.791,0.595) (-0.766,0.626) (-0.740,0.656) (-0.712,0.686) (-0.684,0.714) (-0.654,0.742) (-0.624,0.768) (-0.592,0.793) (-0.560,0.818) (-0.526,0.841) (-0.492,0.862) (-0.457,0.883) (-0.422,0.902) (-0.386,0.919) (-0.349,0.935) (-0.311,0.950) (-0.273,0.963) (-0.235,0.974) (-0.197,0.984) (-0.158,0.992) (-0.118,0.998) (-0.079,1.002) (-0.040,1.005) (0.000,1.006) (0.040,1.005) (0.079,1.002) (0.118,0.998) (0.158,0.992) (0.197,0.984) (0.235,0.974) (0.273,0.963) (0.311,0.950) (0.349,0.935) (0.386,0.919) (0.422,0.902) (0.457,0.883) (0.492,0.862) (0.526,0.841) (0.560,0.818) (0.592,0.793) (0.624,0.768) (0.654,0.742) (0.684,0.714) (0.712,0.686) (0.740,0.656) (0.766,0.626) (0.791,0.595) (0.815,0.563) (0.838,0.531) (0.859,0.498) (0.879,0.465) (0.898,0.431) (0.915,0.396) (0.931,0.361) (0.945,0.326) (0.958,0.291) (0.970,0.255) (0.980,0.219) (0.988,0.183) (0.995,0.146) (1.001,0.110) (1.004,0.073) (1.007,0.037) (1.008,0.000)} -- plot[smooth] coordinates {(1.008,-0.000) (1.007,-0.037) (1.004,-0.073) (1.001,-0.110) (0.995,-0.146) (0.988,-0.183) (0.980,-0.219) (0.970,-0.255) (0.958,-0.291) (0.945,-0.326) (0.931,-0.361) (0.915,-0.396) (0.898,-0.431) (0.879,-0.465) (0.859,-0.498) (0.838,-0.531) (0.815,-0.563) (0.791,-0.595) (0.766,-0.626) (0.740,-0.656) (0.712,-0.686) (0.684,-0.714) (0.654,-0.742) (0.624,-0.768) (0.592,-0.793) (0.560,-0.818) (0.526,-0.841) (0.492,-0.862) (0.457,-0.883) (0.422,-0.902) (0.386,-0.919) (0.349,-0.935) (0.311,-0.950) (0.273,-0.963) (0.235,-0.974) (0.197,-0.984) (0.158,-0.992) (0.118,-0.998) (0.079,-1.002) (0.040,-1.005) (0.000,-1.006) (-0.040,-1.005) (-0.079,-1.002) (-0.118,-0.998) (-0.158,-0.992) (-0.197,-0.984) (-0.235,-0.974) (-0.273,-0.963) (-0.311,-0.950) (-0.349,-0.935) (-0.386,-0.919) (-0.422,-0.902) (-0.457,-0.883) (-0.492,-0.862) (-0.526,-0.841) (-0.560,-0.818) (-0.592,-0.793) (-0.624,-0.768) (-0.654,-0.742) (-0.684,-0.714) (-0.712,-0.686) (-0.740,-0.656) (-0.766,-0.626) (-0.791,-0.595) (-0.815,-0.563) (-0.838,-0.531) (-0.859,-0.498) (-0.879,-0.465) (-0.898,-0.431) (-0.915,-0.396) (-0.931,-0.361) (-0.945,-0.326) (-0.958,-0.291) (-0.970,-0.255) (-0.980,-0.219) (-0.988,-0.183) (-0.995,-0.146) (-1.001,-0.110) (-1.004,-0.073) (-1.007,-0.037) (-1.008,-0.000)} -- cycle;
  \draw[ascblue, thick] plot[smooth] coordinates {(-1.492,0.000) (-1.491,0.045) (-1.488,0.089) (-1.482,0.134) (-1.474,0.179) (-1.464,0.223) (-1.451,0.268) (-1.436,0.313) (-1.419,0.358) (-1.400,0.403) (-1.379,0.448) (-1.355,0.493) (-1.330,0.537) (-1.302,0.582) (-1.273,0.627) (-1.241,0.671) (-1.207,0.715) (-1.172,0.758) (-1.135,0.801) (-1.096,0.843) (-1.055,0.885) (-1.013,0.926) (-0.969,0.966) (-0.924,1.005) (-0.877,1.042) (-0.853,1.060) (-0.829,1.078) (-0.805,1.096) (-0.780,1.113) (-0.755,1.130) (-0.729,1.146) (-0.704,1.162) (-0.678,1.177) (-0.651,1.192) (-0.625,1.206) (-0.598,1.220) (-0.571,1.234) (-0.544,1.247) (-0.517,1.259) (-0.461,1.282) (-0.405,1.302) (-0.348,1.320) (-0.291,1.335) (-0.233,1.348) (-0.175,1.358) (-0.117,1.365) (-0.059,1.369) (0.000,1.371) (0.059,1.369) (0.117,1.365) (0.175,1.358) (0.233,1.348) (0.291,1.335) (0.348,1.320) (0.405,1.302) (0.461,1.282) (0.517,1.259) (0.544,1.247) (0.571,1.234) (0.598,1.220) (0.625,1.206) (0.651,1.192) (0.678,1.177) (0.704,1.162) (0.729,1.146) (0.755,1.130) (0.780,1.113) (0.805,1.096) (0.829,1.078) (0.853,1.060) (0.877,1.042) (0.924,1.005) (0.969,0.966) (1.013,0.926) (1.055,0.885) (1.096,0.843) (1.135,0.801) (1.172,0.758) (1.207,0.715) (1.241,0.671) (1.273,0.627) (1.302,0.582) (1.330,0.537) (1.355,0.493) (1.379,0.448) (1.400,0.403) (1.419,0.358) (1.436,0.313) (1.451,0.268) (1.464,0.223) (1.474,0.179) (1.482,0.134) (1.488,0.089) (1.491,0.045) (1.492,0.000)} -- plot[smooth] coordinates {(1.492,-0.000) (1.491,-0.045) (1.488,-0.089) (1.482,-0.134) (1.474,-0.179) (1.464,-0.223) (1.451,-0.268) (1.436,-0.313) (1.419,-0.358) (1.400,-0.403) (1.379,-0.448) (1.355,-0.493) (1.330,-0.537) (1.302,-0.582) (1.273,-0.627) (1.241,-0.671) (1.207,-0.715) (1.172,-0.758) (1.135,-0.801) (1.096,-0.843) (1.055,-0.885) (1.013,-0.926) (0.969,-0.966) (0.924,-1.005) (0.877,-1.042) (0.853,-1.060) (0.829,-1.078) (0.805,-1.096) (0.780,-1.113) (0.755,-1.130) (0.729,-1.146) (0.704,-1.162) (0.678,-1.177) (0.651,-1.192) (0.625,-1.206) (0.598,-1.220) (0.571,-1.234) (0.544,-1.247) (0.517,-1.259) (0.461,-1.282) (0.405,-1.302) (0.348,-1.320) (0.291,-1.335) (0.233,-1.348) (0.175,-1.358) (0.117,-1.365) (0.059,-1.369) (0.000,-1.371) (-0.059,-1.369) (-0.117,-1.365) (-0.175,-1.358) (-0.233,-1.348) (-0.291,-1.335) (-0.348,-1.320) (-0.405,-1.302) (-0.461,-1.282) (-0.517,-1.259) (-0.544,-1.247) (-0.571,-1.234) (-0.598,-1.220) (-0.625,-1.206) (-0.651,-1.192) (-0.678,-1.177) (-0.704,-1.162) (-0.729,-1.146) (-0.755,-1.130) (-0.780,-1.113) (-0.805,-1.096) (-0.829,-1.078) (-0.853,-1.060) (-0.877,-1.042) (-0.924,-1.005) (-0.969,-0.966) (-1.013,-0.926) (-1.055,-0.885) (-1.096,-0.843) (-1.135,-0.801) (-1.172,-0.758) (-1.207,-0.715) (-1.241,-0.671) (-1.273,-0.627) (-1.302,-0.582) (-1.330,-0.537) (-1.355,-0.493) (-1.379,-0.448) (-1.400,-0.403) (-1.419,-0.358) (-1.436,-0.313) (-1.451,-0.268) (-1.464,-0.223) (-1.474,-0.179) (-1.482,-0.134) (-1.488,-0.089) (-1.491,-0.045) (-1.492,-0.000)} -- cycle;
  \draw[ascwarm, thick] plot[smooth] coordinates {(-2.500,0.000) (-2.469,0.033) (-2.438,0.067) (-2.406,0.100) (-2.375,0.133) (-2.344,0.167) (-2.312,0.200) (-2.281,0.233) (-2.250,0.266) (-2.219,0.299) (-2.188,0.332) (-2.156,0.364) (-2.125,0.397) (-2.094,0.429) (-2.062,0.461) (-2.031,0.493) (-2.000,0.525) (-1.969,0.557) (-1.938,0.588) (-1.906,0.620) (-1.875,0.651) (-1.844,0.681) (-1.812,0.712) (-1.781,0.742) (-1.750,0.772) (-1.719,0.801) (-1.688,0.831) (-1.656,0.860) (-1.625,0.888) (-1.594,0.917) (-1.562,0.944) (-1.531,0.972) (-1.500,0.999) (-1.469,1.026) (-1.438,1.052) (-1.406,1.078) (-1.375,1.104) (-1.344,1.129) (-1.312,1.154) (-1.281,1.178) (-1.250,1.202) (-1.219,1.225) (-1.188,1.248) (-1.156,1.271) (-1.125,1.293) (-1.094,1.314) (-1.062,1.335) (-1.031,1.355) (-1.000,1.375) (-0.969,1.395) (-0.938,1.413) (-0.906,1.432) (-0.875,1.449) (-0.844,1.467) (-0.812,1.483) (-0.781,1.499) (-0.750,1.515) (-0.719,1.530) (-0.688,1.544) (-0.656,1.558) (-0.625,1.571) (-0.594,1.583) (-0.562,1.595) (-0.531,1.606) (-0.500,1.617) (-0.469,1.627) (-0.438,1.636) (-0.406,1.645) (-0.375,1.653) (-0.344,1.661) (-0.312,1.667) (-0.281,1.674) (-0.250,1.679) (-0.219,1.684) (-0.188,1.688) (-0.156,1.692) (-0.125,1.695) (-0.094,1.697) (-0.062,1.699) (-0.031,1.700) (0.000,1.700) (0.031,1.700) (0.062,1.699) (0.094,1.697) (0.125,1.695) (0.156,1.692) (0.188,1.688) (0.219,1.684) (0.250,1.679) (0.281,1.674) (0.312,1.667) (0.344,1.661) (0.375,1.653) (0.406,1.645) (0.438,1.636) (0.469,1.627) (0.500,1.617) (0.531,1.606) (0.562,1.595) (0.594,1.583) (0.625,1.571) (0.656,1.558) (0.687,1.544) (0.719,1.530) (0.750,1.515) (0.781,1.499) (0.813,1.483) (0.844,1.467) (0.875,1.449) (0.906,1.432) (0.937,1.413) (0.969,1.395) (1.000,1.375) (1.031,1.355) (1.063,1.335) (1.094,1.314) (1.125,1.293) (1.156,1.271) (1.187,1.248) (1.219,1.225) (1.250,1.202) (1.281,1.178) (1.313,1.154) (1.344,1.129) (1.375,1.104) (1.406,1.078) (1.437,1.052) (1.469,1.026) (1.500,0.999) (1.531,0.972) (1.563,0.944) (1.594,0.917) (1.625,0.888) (1.656,0.860) (1.687,0.831) (1.719,0.801) (1.750,0.772) (1.781,0.742) (1.813,0.712) (1.844,0.681) (1.875,0.651) (1.906,0.620) (1.937,0.588) (1.969,0.557) (2.000,0.525) (2.031,0.493) (2.063,0.461) (2.094,0.429) (2.125,0.397) (2.156,0.364) (2.187,0.332) (2.219,0.299) (2.250,0.266) (2.281,0.233) (2.313,0.200) (2.344,0.167) (2.375,0.133) (2.406,0.100) (2.438,0.067) (2.469,0.033) (2.500,0.000)};
  \draw[ascwarm, thick] plot[smooth] coordinates {(-2.500,-0.000) (-2.469,-0.033) (-2.438,-0.067) (-2.406,-0.100) (-2.375,-0.133) (-2.344,-0.167) (-2.312,-0.200) (-2.281,-0.233) (-2.250,-0.266) (-2.219,-0.299) (-2.188,-0.332) (-2.156,-0.364) (-2.125,-0.397) (-2.094,-0.429) (-2.062,-0.461) (-2.031,-0.493) (-2.000,-0.525) (-1.969,-0.557) (-1.938,-0.588) (-1.906,-0.620) (-1.875,-0.651) (-1.844,-0.681) (-1.812,-0.712) (-1.781,-0.742) (-1.750,-0.772) (-1.719,-0.801) (-1.688,-0.831) (-1.656,-0.860) (-1.625,-0.888) (-1.594,-0.917) (-1.562,-0.944) (-1.531,-0.972) (-1.500,-0.999) (-1.469,-1.026) (-1.438,-1.052) (-1.406,-1.078) (-1.375,-1.104) (-1.344,-1.129) (-1.312,-1.154) (-1.281,-1.178) (-1.250,-1.202) (-1.219,-1.225) (-1.188,-1.248) (-1.156,-1.271) (-1.125,-1.293) (-1.094,-1.314) (-1.062,-1.335) (-1.031,-1.355) (-1.000,-1.375) (-0.969,-1.395) (-0.938,-1.413) (-0.906,-1.432) (-0.875,-1.449) (-0.844,-1.467) (-0.812,-1.483) (-0.781,-1.499) (-0.750,-1.515) (-0.719,-1.530) (-0.688,-1.544) (-0.656,-1.558) (-0.625,-1.571) (-0.594,-1.583) (-0.562,-1.595) (-0.531,-1.606) (-0.500,-1.617) (-0.469,-1.627) (-0.438,-1.636) (-0.406,-1.645) (-0.375,-1.653) (-0.344,-1.661) (-0.312,-1.667) (-0.281,-1.674) (-0.250,-1.679) (-0.219,-1.684) (-0.188,-1.688) (-0.156,-1.692) (-0.125,-1.695) (-0.094,-1.697) (-0.062,-1.699) (-0.031,-1.700) (0.000,-1.700) (0.031,-1.700) (0.062,-1.699) (0.094,-1.697) (0.125,-1.695) (0.156,-1.692) (0.188,-1.688) (0.219,-1.684) (0.250,-1.679) (0.281,-1.674) (0.312,-1.667) (0.344,-1.661) (0.375,-1.653) (0.406,-1.645) (0.438,-1.636) (0.469,-1.627) (0.500,-1.617) (0.531,-1.606) (0.562,-1.595) (0.594,-1.583) (0.625,-1.571) (0.656,-1.558) (0.687,-1.544) (0.719,-1.530) (0.750,-1.515) (0.781,-1.499) (0.813,-1.483) (0.844,-1.467) (0.875,-1.449) (0.906,-1.432) (0.937,-1.413) (0.969,-1.395) (1.000,-1.375) (1.031,-1.355) (1.063,-1.335) (1.094,-1.314) (1.125,-1.293) (1.156,-1.271) (1.187,-1.248) (1.219,-1.225) (1.250,-1.202) (1.281,-1.178) (1.313,-1.154) (1.344,-1.129) (1.375,-1.104) (1.406,-1.078) (1.437,-1.052) (1.469,-1.026) (1.500,-0.999) (1.531,-0.972) (1.563,-0.944) (1.594,-0.917) (1.625,-0.888) (1.656,-0.860) (1.687,-0.831) (1.719,-0.801) (1.750,-0.772) (1.781,-0.742) (1.813,-0.712) (1.844,-0.681) (1.875,-0.651) (1.906,-0.620) (1.937,-0.588) (1.969,-0.557) (2.000,-0.525) (2.031,-0.493) (2.063,-0.461) (2.094,-0.429) (2.125,-0.397) (2.156,-0.364) (2.187,-0.332) (2.219,-0.299) (2.250,-0.266) (2.281,-0.233) (2.313,-0.200) (2.344,-0.167) (2.375,-0.133) (2.406,-0.100) (2.438,-0.067) (2.469,-0.033) (2.500,-0.000)};
  \draw[ascpurple, thick] plot[smooth] coordinates {(-2.500,1.006) (-2.469,1.006) (-2.438,1.008) (-2.406,1.011) (-2.375,1.015) (-2.344,1.019) (-2.312,1.025) (-2.281,1.032) (-2.250,1.040) (-2.219,1.049) (-2.188,1.059) (-2.156,1.070) (-2.125,1.081) (-2.094,1.094) (-2.062,1.107) (-2.031,1.120) (-2.000,1.135) (-1.969,1.150) (-1.938,1.165) (-1.906,1.181) (-1.875,1.198) (-1.844,1.215) (-1.812,1.232) (-1.781,1.250) (-1.750,1.268) (-1.719,1.286) (-1.688,1.304) (-1.656,1.323) (-1.625,1.342) (-1.594,1.361) (-1.562,1.380) (-1.531,1.399) (-1.500,1.418) (-1.469,1.437) (-1.438,1.456) (-1.406,1.475) (-1.375,1.493) (-1.344,1.512) (-1.312,1.531) (-1.281,1.549) (-1.250,1.567) (-1.219,1.585) (-1.188,1.603) (-1.156,1.621) (-1.125,1.638) (-1.094,1.655) (-1.062,1.671) (-1.031,1.688) (-1.000,1.704) (-0.969,1.719) (-0.938,1.735) (-0.906,1.750) (-0.875,1.764) (-0.844,1.778) (-0.812,1.792) (-0.781,1.805) (-0.750,1.818) (-0.719,1.831) (-0.688,1.843) (-0.656,1.854) (-0.625,1.865) (-0.594,1.876) (-0.562,1.886) (-0.531,1.895) (-0.500,1.904) (-0.469,1.913) (-0.438,1.921) (-0.406,1.928) (-0.375,1.935) (-0.344,1.941) (-0.312,1.947) (-0.281,1.952) (-0.250,1.957) (-0.219,1.961) (-0.188,1.965) (-0.156,1.968) (-0.125,1.971) (-0.094,1.973) (-0.062,1.974) (-0.031,1.975) (0.000,1.975) (0.031,1.975) (0.062,1.974) (0.094,1.973) (0.125,1.971) (0.156,1.968) (0.188,1.965) (0.219,1.961) (0.250,1.957) (0.281,1.952) (0.312,1.947) (0.344,1.941) (0.375,1.935) (0.406,1.928) (0.438,1.921) (0.469,1.913) (0.500,1.904) (0.531,1.895) (0.562,1.886) (0.594,1.876) (0.625,1.865) (0.656,1.854) (0.687,1.843) (0.719,1.831) (0.750,1.818) (0.781,1.805) (0.813,1.792) (0.844,1.778) (0.875,1.764) (0.906,1.750) (0.937,1.735) (0.969,1.719) (1.000,1.704) (1.031,1.688) (1.063,1.671) (1.094,1.655) (1.125,1.638) (1.156,1.621) (1.187,1.603) (1.219,1.585) (1.250,1.567) (1.281,1.549) (1.313,1.531) (1.344,1.512) (1.375,1.493) (1.406,1.475) (1.437,1.456) (1.469,1.437) (1.500,1.418) (1.531,1.399) (1.563,1.380) (1.594,1.361) (1.625,1.342) (1.656,1.323) (1.687,1.304) (1.719,1.286) (1.750,1.268) (1.781,1.250) (1.813,1.232) (1.844,1.215) (1.875,1.198) (1.906,1.181) (1.937,1.165) (1.969,1.150) (2.000,1.135) (2.031,1.120) (2.063,1.107) (2.094,1.094) (2.125,1.081) (2.156,1.070) (2.187,1.059) (2.219,1.049) (2.250,1.040) (2.281,1.032) (2.313,1.025) (2.344,1.019) (2.375,1.015) (2.406,1.011) (2.438,1.008) (2.469,1.006) (2.500,1.006)};
  \draw[ascpurple, thick] plot[smooth] coordinates {(-2.500,-1.006) (-2.469,-1.006) (-2.438,-1.008) (-2.406,-1.011) (-2.375,-1.015) (-2.344,-1.019) (-2.312,-1.025) (-2.281,-1.032) (-2.250,-1.040) (-2.219,-1.049) (-2.188,-1.059) (-2.156,-1.070) (-2.125,-1.081) (-2.094,-1.094) (-2.062,-1.107) (-2.031,-1.120) (-2.000,-1.135) (-1.969,-1.150) (-1.938,-1.165) (-1.906,-1.181) (-1.875,-1.198) (-1.844,-1.215) (-1.812,-1.232) (-1.781,-1.250) (-1.750,-1.268) (-1.719,-1.286) (-1.688,-1.304) (-1.656,-1.323) (-1.625,-1.342) (-1.594,-1.361) (-1.562,-1.380) (-1.531,-1.399) (-1.500,-1.418) (-1.469,-1.437) (-1.438,-1.456) (-1.406,-1.475) (-1.375,-1.493) (-1.344,-1.512) (-1.312,-1.531) (-1.281,-1.549) (-1.250,-1.567) (-1.219,-1.585) (-1.188,-1.603) (-1.156,-1.621) (-1.125,-1.638) (-1.094,-1.655) (-1.062,-1.671) (-1.031,-1.688) (-1.000,-1.704) (-0.969,-1.719) (-0.938,-1.735) (-0.906,-1.750) (-0.875,-1.764) (-0.844,-1.778) (-0.812,-1.792) (-0.781,-1.805) (-0.750,-1.818) (-0.719,-1.831) (-0.688,-1.843) (-0.656,-1.854) (-0.625,-1.865) (-0.594,-1.876) (-0.562,-1.886) (-0.531,-1.895) (-0.500,-1.904) (-0.469,-1.913) (-0.438,-1.921) (-0.406,-1.928) (-0.375,-1.935) (-0.344,-1.941) (-0.312,-1.947) (-0.281,-1.952) (-0.250,-1.957) (-0.219,-1.961) (-0.188,-1.965) (-0.156,-1.968) (-0.125,-1.971) (-0.094,-1.973) (-0.062,-1.974) (-0.031,-1.975) (0.000,-1.975) (0.031,-1.975) (0.062,-1.974) (0.094,-1.973) (0.125,-1.971) (0.156,-1.968) (0.188,-1.965) (0.219,-1.961) (0.250,-1.957) (0.281,-1.952) (0.312,-1.947) (0.344,-1.941) (0.375,-1.935) (0.406,-1.928) (0.438,-1.921) (0.469,-1.913) (0.500,-1.904) (0.531,-1.895) (0.562,-1.886) (0.594,-1.876) (0.625,-1.865) (0.656,-1.854) (0.687,-1.843) (0.719,-1.831) (0.750,-1.818) (0.781,-1.805) (0.813,-1.792) (0.844,-1.778) (0.875,-1.764) (0.906,-1.750) (0.937,-1.735) (0.969,-1.719) (1.000,-1.704) (1.031,-1.688) (1.063,-1.671) (1.094,-1.655) (1.125,-1.638) (1.156,-1.621) (1.187,-1.603) (1.219,-1.585) (1.250,-1.567) (1.281,-1.549) (1.313,-1.531) (1.344,-1.512) (1.375,-1.493) (1.406,-1.475) (1.437,-1.456) (1.469,-1.437) (1.500,-1.418) (1.531,-1.399) (1.563,-1.380) (1.594,-1.361) (1.625,-1.342) (1.656,-1.323) (1.687,-1.304) (1.719,-1.286) (1.750,-1.268) (1.781,-1.250) (1.813,-1.232) (1.844,-1.215) (1.875,-1.198) (1.906,-1.181) (1.937,-1.165) (1.969,-1.150) (2.000,-1.135) (2.031,-1.120) (2.063,-1.107) (2.094,-1.094) (2.125,-1.081) (2.156,-1.070) (2.187,-1.059) (2.219,-1.049) (2.250,-1.040) (2.281,-1.032) (2.313,-1.025) (2.344,-1.019) (2.375,-1.015) (2.406,-1.011) (2.438,-1.008) (2.469,-1.006) (2.500,-1.006)};
  \draw[ascpurple, thick] plot[smooth] coordinates {(-2.500,1.521) (-2.469,1.521) (-2.438,1.522) (-2.406,1.524) (-2.375,1.526) (-2.344,1.530) (-2.312,1.534) (-2.281,1.538) (-2.250,1.544) (-2.219,1.550) (-2.188,1.556) (-2.156,1.564) (-2.125,1.571) (-2.094,1.580) (-2.062,1.589) (-2.031,1.599) (-2.000,1.609) (-1.969,1.619) (-1.938,1.630) (-1.906,1.642) (-1.875,1.654) (-1.844,1.666) (-1.812,1.679) (-1.781,1.692) (-1.750,1.705) (-1.719,1.719) (-1.688,1.733) (-1.656,1.747) (-1.625,1.761) (-1.594,1.775) (-1.562,1.790) (-1.531,1.805) (-1.500,1.819) (-1.469,1.834) (-1.438,1.849) (-1.406,1.864) (-1.375,1.879) (-1.344,1.894) (-1.312,1.909) (-1.281,1.924) (-1.250,1.938) (-1.219,1.953) (-1.188,1.967) (-1.156,1.982) (-1.125,1.996) (-1.094,2.010) (-1.062,2.023) (-1.031,2.037) (-1.000,2.050) (-0.969,2.063) (-0.938,2.076) (-0.906,2.089) (-0.875,2.101) (-0.844,2.113) (-0.812,2.124) (-0.781,2.135) (-0.750,2.146) (-0.719,2.157) (-0.688,2.167) (-0.656,2.177) (-0.625,2.186) (-0.594,2.195) (-0.562,2.204) (-0.531,2.212) (-0.500,2.219) (-0.469,2.227) (-0.438,2.234) (-0.406,2.240) (-0.375,2.246) (-0.344,2.252) (-0.312,2.257) (-0.281,2.261) (-0.250,2.265) (-0.219,2.269) (-0.188,2.272) (-0.156,2.275) (-0.125,2.277) (-0.094,2.279) (-0.062,2.280) (-0.031,2.281) (0.000,2.281) (0.031,2.281) (0.062,2.280) (0.094,2.279) (0.125,2.277) (0.156,2.275) (0.188,2.272) (0.219,2.269) (0.250,2.265) (0.281,2.261) (0.312,2.257) (0.344,2.252) (0.375,2.246) (0.406,2.240) (0.438,2.234) (0.469,2.227) (0.500,2.219) (0.531,2.212) (0.562,2.204) (0.594,2.195) (0.625,2.186) (0.656,2.177) (0.687,2.167) (0.719,2.157) (0.750,2.146) (0.781,2.135) (0.813,2.124) (0.844,2.113) (0.875,2.101) (0.906,2.089) (0.937,2.076) (0.969,2.063) (1.000,2.050) (1.031,2.037) (1.063,2.023) (1.094,2.010) (1.125,1.996) (1.156,1.982) (1.187,1.967) (1.219,1.953) (1.250,1.938) (1.281,1.924) (1.313,1.909) (1.344,1.894) (1.375,1.879) (1.406,1.864) (1.437,1.849) (1.469,1.834) (1.500,1.819) (1.531,1.805) (1.563,1.790) (1.594,1.775) (1.625,1.761) (1.656,1.747) (1.687,1.733) (1.719,1.719) (1.750,1.705) (1.781,1.692) (1.813,1.679) (1.844,1.666) (1.875,1.654) (1.906,1.642) (1.937,1.630) (1.969,1.619) (2.000,1.609) (2.031,1.599) (2.063,1.589) (2.094,1.580) (2.125,1.571) (2.156,1.564) (2.187,1.556) (2.219,1.550) (2.250,1.544) (2.281,1.538) (2.313,1.534) (2.344,1.530) (2.375,1.526) (2.406,1.524) (2.438,1.522) (2.469,1.521) (2.500,1.521)};
  \draw[ascpurple, thick] plot[smooth] coordinates {(-2.500,-1.521) (-2.469,-1.521) (-2.438,-1.522) (-2.406,-1.524) (-2.375,-1.526) (-2.344,-1.530) (-2.312,-1.534) (-2.281,-1.538) (-2.250,-1.544) (-2.219,-1.550) (-2.188,-1.556) (-2.156,-1.564) (-2.125,-1.571) (-2.094,-1.580) (-2.062,-1.589) (-2.031,-1.599) (-2.000,-1.609) (-1.969,-1.619) (-1.938,-1.630) (-1.906,-1.642) (-1.875,-1.654) (-1.844,-1.666) (-1.812,-1.679) (-1.781,-1.692) (-1.750,-1.705) (-1.719,-1.719) (-1.688,-1.733) (-1.656,-1.747) (-1.625,-1.761) (-1.594,-1.775) (-1.562,-1.790) (-1.531,-1.805) (-1.500,-1.819) (-1.469,-1.834) (-1.438,-1.849) (-1.406,-1.864) (-1.375,-1.879) (-1.344,-1.894) (-1.312,-1.909) (-1.281,-1.924) (-1.250,-1.938) (-1.219,-1.953) (-1.188,-1.967) (-1.156,-1.982) (-1.125,-1.996) (-1.094,-2.010) (-1.062,-2.023) (-1.031,-2.037) (-1.000,-2.050) (-0.969,-2.063) (-0.938,-2.076) (-0.906,-2.089) (-0.875,-2.101) (-0.844,-2.113) (-0.812,-2.124) (-0.781,-2.135) (-0.750,-2.146) (-0.719,-2.157) (-0.688,-2.167) (-0.656,-2.177) (-0.625,-2.186) (-0.594,-2.195) (-0.562,-2.204) (-0.531,-2.212) (-0.500,-2.219) (-0.469,-2.227) (-0.438,-2.234) (-0.406,-2.240) (-0.375,-2.246) (-0.344,-2.252) (-0.312,-2.257) (-0.281,-2.261) (-0.250,-2.265) (-0.219,-2.269) (-0.188,-2.272) (-0.156,-2.275) (-0.125,-2.277) (-0.094,-2.279) (-0.062,-2.280) (-0.031,-2.281) (0.000,-2.281) (0.031,-2.281) (0.062,-2.280) (0.094,-2.279) (0.125,-2.277) (0.156,-2.275) (0.188,-2.272) (0.219,-2.269) (0.250,-2.265) (0.281,-2.261) (0.312,-2.257) (0.344,-2.252) (0.375,-2.246) (0.406,-2.240) (0.438,-2.234) (0.469,-2.227) (0.500,-2.219) (0.531,-2.212) (0.562,-2.204) (0.594,-2.195) (0.625,-2.186) (0.656,-2.177) (0.687,-2.167) (0.719,-2.157) (0.750,-2.146) (0.781,-2.135) (0.813,-2.124) (0.844,-2.113) (0.875,-2.101) (0.906,-2.089) (0.937,-2.076) (0.969,-2.063) (1.000,-2.050) (1.031,-2.037) (1.063,-2.023) (1.094,-2.010) (1.125,-1.996) (1.156,-1.982) (1.187,-1.967) (1.219,-1.953) (1.250,-1.938) (1.281,-1.924) (1.313,-1.909) (1.344,-1.894) (1.375,-1.879) (1.406,-1.864) (1.437,-1.849) (1.469,-1.834) (1.500,-1.819) (1.531,-1.805) (1.563,-1.790) (1.594,-1.775) (1.625,-1.761) (1.656,-1.747) (1.687,-1.733) (1.719,-1.719) (1.750,-1.705) (1.781,-1.692) (1.813,-1.679) (1.844,-1.666) (1.875,-1.654) (1.906,-1.642) (1.937,-1.630) (1.969,-1.619) (2.000,-1.609) (2.031,-1.599) (2.063,-1.589) (2.094,-1.580) (2.125,-1.571) (2.156,-1.564) (2.187,-1.556) (2.219,-1.550) (2.250,-1.544) (2.281,-1.538) (2.313,-1.534) (2.344,-1.530) (2.375,-1.526) (2.406,-1.524) (2.438,-1.522) (2.469,-1.521) (2.500,-1.521)};
  \draw[ascblue, -{Stealth[length=6pt]}] (0.02,1.371) -- (0.25,1.371);
  \draw[ascpurple, -{Stealth[length=6pt]}] (1.2,1.596) -- (1.5,1.418);
  \node[ascwarm, font=\small, anchor=west] at (2.55,0.55) {separatrix};
\end{tikzpicture}

{\small The phase space of a pendulum. Each curve is a path of constant energy. Low energies give closed loops (blue: swinging back and forth); high energies give wavy lines (purple: spinning over the top); the separatrix divides the two.}
\end{minipage}
\end{center}

\begin{theorem}[Liouville's theorem]
The Hamiltonian flow preserves volume in phase space. A blob of initial conditions may stretch and fold, but its volume never changes.
\end{theorem}

Liouville's theorem has big consequences. A Hamiltonian system cannot have an attractor --- a region that sucks in all nearby trajectories --- because that would shrink phase-space volume; attractors need friction (Chapter~\ref{ch:chaos}). And it is the reason statistical mechanics can assign equal probability to equal volumes of phase space (Chapter~\ref{ch:statmech}): the dynamics never makes some regions ``more likely'' than others.

\section{Poisson brackets: the bridge to quantum mechanics}

For any two functions $f(q,p)$ and $g(q,p)$ on phase space, define the \textbf{Poisson bracket}
\[
  \{f,g\} = \sum_i\left(\pdv{f}{q_i}\pdv{g}{p_i} - \pdv{f}{p_i}\pdv{g}{q_i}\right).
\]
Then the time evolution of \emph{any} quantity is
\[
  \dv{f}{t} = \{f,H\},
\]
so $f$ is conserved exactly when its bracket with $H$ vanishes. The basic brackets are $\{q_i,p_j\} = \delta_{ij}$, and for angular momentum $\{L_x, L_y\} = L_z$ and cyclic permutations.

Now the astonishing part. In 1925 Dirac noticed that quantum mechanics has exactly the same structure, with the Poisson bracket replaced by a commutator of operators:
\[
  \{f,g\} \quad\longrightarrow\quad \frac{1}{i\hbar}[\hat f,\hat g], \qquad\text{so}\qquad \{x,p\}=1 \;\longrightarrow\; [\hat x,\hat p] = i\hbar .
\]
Every classical structure in this section survives quantization almost unchanged. That is why we learn Hamiltonian mechanics: it is the classical shadow of quantum mechanics (Chapter~\ref{ch:qm}).

\section{Small oscillations and normal modes}

Near a stable equilibrium, any system with $n$ degrees of freedom has, to second order,
\[
  L = \tfrac12\,\dot{\vb q}^{T} M\,\dot{\vb q} - \tfrac12\,\vb q^{T}K\,\vb q,
\]
with $M$ and $K$ symmetric matrices. The equations of motion $M\ddot{\vb q} = -K\vb q$ are solved by $\vb q = \vb a\,e^{i\omega t}$ if
\[
  \det(K - \omega^2 M) = 0 .
\]
The $n$ solutions $\omega_k$ are the \emph{normal mode frequencies}, and in the coordinates of the eigenvectors the system falls apart into $n$ independent harmonic oscillators. Every vibrating system --- a molecule, a bridge, a crystal, a drum --- is a collection of independent oscillators in the right coordinates. This is the spherical cow of Chapter~\ref{ch:sho}, multiplied.

\section{Adiabatic invariants}

Suppose a parameter of a system changes slowly compared to its natural period: a pendulum whose string is slowly shortened, say. The energy is not conserved --- work is being done --- but something else nearly is. For any periodic motion, the \textbf{action variable}
\[
  J = \oint p\,\dd q,
\]
the area enclosed by the orbit in phase space, is an \emph{adiabatic invariant}: it changes by an amount that vanishes faster than any power of the slowness. For an oscillator, the ellipse has area $2\pi E/\omega$, so $E/\omega$ stays constant: shorten a pendulum slowly and its energy rises in proportion to its frequency.

This had a remarkable historical role. In early quantum theory, the quantities that were quantized ($E = n\hbar\omega$ for oscillators, $J = nh$ in general) had to be adiabatic invariants, or else slowly changing the parameters could carry a system between allowed states into forbidden ones.

\section{Why the action? A glimpse of the path integral}

Why should nature extremize anything? The deepest answer comes from quantum mechanics. In Feynman's formulation, a quantum particle does not take one path: it takes \emph{all} of them, and each path contributes a complex number $e^{iS/\hbar}$ to the probability amplitude. For macroscopic objects, $S\gg\hbar$, and the phase $S/\hbar$ swings wildly between neighbouring paths, so their contributions cancel --- except near paths where $S$ is stationary, where neighbouring paths have nearly the same phase and add up. The classical path is the one that survives the cancellation.

\begin{keyidea}
Classical mechanics is the $\hbar\to0$ limit of quantum mechanics, and the principle of stationary action is what that limit looks like. Nature is not lazy; it is democratic, and the classical path wins by constructive interference.
\end{keyidea}

\section{Problems}

\begin{problems}
\prob{1} Write down the Lagrangian of an Atwood machine (two masses on a pulley) using a single coordinate, derive its acceleration, and confirm it agrees with Newton's laws. Notice that you never needed the tension.

\prob{1}\core For a particle moving in a central potential in three dimensions, use spherical coordinates to show that $L = \tfrac12 m(\dot r^2 + r^2\dot\theta^2 + r^2\sin^2\theta\,\dot\phi^2) - V(r)$, and identify which momentum is immediately conserved. Why are all three components of angular momentum conserved, even though only one is manifest?

\prob{1} Write the Hamiltonian of the harmonic oscillator, derive Hamilton's equations, and verify that the phase-space trajectories are ellipses traversed clockwise. What is the area of the ellipse in terms of $E$ and $\omega$?

\prob{2} Redo the bead on a rotating hoop (Problem~\ref{prob:bead-hoop}) using a Lagrangian in the rotating frame. The Lagrangian has no explicit time dependence, so $H$ is conserved; show that $H$ is \emph{not} the bead's energy, and explain why the energy is not conserved even though $H$ is.

\prob{2}\core \emph{(The double pendulum.)} Two equal masses $m$ hang in series from equal massless rods of length $\ell$. Write the Lagrangian, linearize for small angles, and show that the normal mode frequencies are $\omega^2 = (2\pm\sqrt2)\,g/\ell$. Describe each mode. (For large angles, the double pendulum is chaotic; see Chapter~\ref{ch:chaos}.)

\prob{2}\core Compute the Poisson brackets $\{L_x, L_y\}$, $\{L_z, L^2\}$, and $\{x, L_z\}$, where $\vb L = \vb r\times\vb p$. These become the angular momentum commutation relations of quantum mechanics, and they are why a quantum particle cannot have a definite $L_x$ and $L_y$ at the same time.

\prob{2} A spherical pendulum (a mass on a rod free to swing in any direction) has two conserved quantities. Identify them, reduce the motion to an effective one-dimensional problem in the polar angle $\theta$, and find the angular frequency of a steady conical orbit at angle $\theta_0$.

\prob{2} \emph{(Fermat implies Snell.)} Light travels at speed $c/n_1$ in one medium and $c/n_2$ in another, separated by a flat interface. Minimize the travel time between a point in each medium to derive Snell's law, $n_1\sin\theta_1 = n_2\sin\theta_2$.

\prob{3}\core \emph{(The brachistochrone.)} Find the curve $y(x)$ down which a bead, starting from rest, slides frictionlessly from one point to a lower one in the least time. Show that the time is $\int\sqrt{(1+y'^2)/(2gy)}\,\dd x$ (with $y$ measured downward), use the fact that the integrand does not depend on $x$ to obtain a first integral, and show that the solution is a cycloid. Compare with Problem~\ref{prob:tautochrone}.

\prob{3}\label{prob:kapitza} \emph{(The Kapitza pendulum.)} The pivot of a pendulum is shaken vertically as $a\cos\Omega t$, with $a\ll\ell$ and $\Omega$ much larger than the pendulum's natural frequency. Separating the motion into a slow part and a fast jitter and averaging over the fast jitter, show that the slow motion feels an effective potential
\[
  U_{\text{eff}}(\theta) = -mg\ell\cos\theta + \frac{m a^2\Omega^2}{4}\sin^2\theta,
\]
and conclude that the \emph{inverted} position $\theta=\pi$ becomes stable when $a^2\Omega^2 > 2g\ell$. A broomstick balanced on a vibrating saw stays upright.
\hint{Write $\theta = \bar\theta + \xi$ with $\xi$ small and fast; solve for $\xi$ to leading order, then average the force over one fast period.}

\prob{3} \emph{(An adiabatic piston.)} A ball bounces elastically back and forth between two walls a distance $\ell$ apart. One wall is moved slowly, at a speed much smaller than the ball's. Show that $v\ell$ is constant, so the ball's energy scales as $1/\ell^2$. Interpreting a gas as many such balls, show that slowly compressing a one-dimensional gas gives $P\ell^3 = $ constant, and compare with the adiabatic law $PV^\gamma=$ constant of Chapter~\ref{ch:thermo}. What is $\gamma$ for a one-dimensional gas?
\end{problems}

\begin{reading}
\rd{John R.\ Taylor, \emph{Classical Mechanics}}{Chapters~6, 7, 11, and 13 give the gentlest complete introduction to Lagrangians, Hamiltonians, and normal modes.}
\rd{\emph{The Feynman Lectures on Physics}, Volume II, Chapter~19}{``The Principle of Least Action,'' a famous lecture that explains why the action matters better than any textbook.}
\rd{L.\ D.\ Landau and E.\ M.\ Lifshitz, \emph{Mechanics}}{Short, dense, and magnificent. Derives all of mechanics from symmetry and the action. Read it after a first course.}
\rd{Helliwell and Sahakian, \emph{Modern Classical Mechanics}}{A graduate-level treatment with modern examples and close attention to the bridge to quantum mechanics.}
\end{reading}
